\begin{thebibliography}{99}
\raggedright
\bibitem{makhlin} Y. Makhlin, Nonlocal properties of two-qubit gates and
mixed states and optimization of quantum computations, Quantum Information
Processing 1, 243--252 (2002), \href{https://arxiv.org/abs/quant-ph/0002045}{arXiv:quant-ph/0002045}.
\bibitem{shende} V. V. Shende, S. S. Bullock, and I. L. Markov, Recognizing
small-circuit structure in two-qubit operators and timing Hamiltonians to
compute controlled-NOT gates, Physical Review A 70, 012310 (2004),
\href{https://doi.org/10.1103/PhysRevA.70.012310}{doi:10.1103/PhysRevA.70.012310}.
\bibitem{bgate} J. Zhang, J. Vala, S. Sastry, and K. B. Whaley, Minimum
construction of two-qubit quantum operations, Physical Review Letters 93,
020502 (2004), \href{https://arxiv.org/abs/quant-ph/0312193}{arXiv:quant-ph/0312193}.
\bibitem{sqisw} C. Huang et al., Quantum instruction set design for
performance, Physical Review Letters 130, 070601 (2023),
\href{https://arxiv.org/abs/2105.06074}{arXiv:2105.06074}.
\bibitem{quartz} M. Xu, Z. Li, O. Padon, S. Lin, J. Pointing, A. Hirth,
H. Ma, J. Palsberg, A. Aiken, U. A. Acar, and Z. Jia, Quartz:
Superoptimization of quantum circuits, Proceedings of PLDI (2022),
\href{https://doi.org/10.1145/3519939.3523433}{doi:10.1145/3519939.3523433}.
\bibitem{queso} A. Xu, A. Molavi, L. Pick, S. Tannu, and A. Albarghouthi,
Synthesizing quantum-circuit optimizers, Proceedings of the ACM on
Programming Languages 7, PLDI, Article 140 (2023),
\href{https://doi.org/10.1145/3591254}{doi:10.1145/3591254}.
\bibitem{universalqcompiler} R. Iten, O. Reardon-Smith, E. Malvetti,
L. Mondada, G. Pauvert, E. Redmond, R. S. Kohli, and R. Colbeck,
Introduction to UniversalQCompiler,
\href{https://arxiv.org/abs/1904.01072}{arXiv:1904.01072} (2019; revised 2021).
\bibitem{qiskitparameter} IBM Quantum, Qiskit Parameter documentation,
\url{https://quantum.cloud.ibm.com/docs/en/api/qiskit/qiskit.circuit.Parameter},
accessed 7 November 2026; executed versions are in the experiment manifest.
\bibitem{virtualz} D. C. McKay, C. J. Wood, S. Sheldon, J. M. Chow, and
J. M. Gambetta, Efficient Z gates for quantum computing, Physical Review A
96, 022330 (2017), \href{https://arxiv.org/abs/1612.00858}{arXiv:1612.00858}.
\bibitem{partition} G. Morse, M. Al Farizi, J. N\'adori, L. Hajas,
T. Kozsik, and P. Rakyta, Entanglement-aware quantum circuit optimization
through globally optimal partitioning,
\href{https://arxiv.org/abs/2611.03845}{arXiv:2611.03845} (2026).

\bibitem{nielsen} M. A. Nielsen, C. M. Dawson, J. L. Dodd, A. Gilchrist,
D. Mortimer, T. J. Osborne, M. J. Bremner, A. W. Harrow, and A. Hines,
Quantum dynamics as a physical resource, Physical Review A 67, 052301 (2003),
\href{https://arxiv.org/abs/quant-ph/0208077}{arXiv:quant-ph/0208077}.
\bibitem{yuduan} N. Yu, R. Duan, and M. Ying, Five two-qubit gates are
necessary for implementing Toffoli gate, Physical Review A 88, 010304(R)
(2013), \href{https://arxiv.org/abs/1301.3372}{arXiv:1301.3372}.
\bibitem{yuying} N. Yu and M. Ying, Optimal simulation of Deutsch gates and
the Fredkin gate, Physical Review A 91, 032302 (2015),
\href{https://doi.org/10.1103/PhysRevA.91.032302}{doi:10.1103/PhysRevA.91.032302}.
The controlled-product statement is Proposition 2.
\bibitem{toffolicnot} V. V. Shende and I. L. Markov, On the CNOT-cost of
TOFFOLI gates, Quantum Information and Computation 9, 461--486 (2009),
\href{https://arxiv.org/abs/0803.2316}{arXiv:0803.2316}.
\bibitem{matsumoto} K. Matsumoto and K. Amano, Representation of quantum
circuits with Clifford and $\pi/8$ gates (2008),
\href{https://arxiv.org/abs/0806.3834}{arXiv:0806.3834}.
\bibitem{kmm} V. Kliuchnikov, D. Maslov, and M. Mosca, Fast and efficient
exact synthesis of single-qubit unitaries generated by Clifford and T gates,
Quantum Information and Computation 13, 607--630 (2013),
\href{https://arxiv.org/abs/1206.5236}{arXiv:1206.5236}.
\bibitem{giles} B. Giles and P. Selinger, Exact synthesis of multiqubit
Clifford+T circuits, Physical Review A 87, 032332 (2013),
\href{https://arxiv.org/abs/1212.0506}{arXiv:1212.0506}.
\bibitem{gilesma} B. Giles and P. Selinger, Remarks on Matsumoto and Amano's
normal form for single-qubit Clifford+T operators (2013),
\href{https://arxiv.org/abs/1312.6584}{arXiv:1312.6584}.
\bibitem{maslov} D. Maslov, Advantages of using relative-phase Toffoli gates
with an application to multiple control Toffoli optimization, Physical Review
A 93, 022311 (2016), \href{https://arxiv.org/abs/1508.03273}{arXiv:1508.03273}.
\bibitem{qiskitsynthesis} IBM Quantum, Qiskit synthesis API,
\url{https://quantum.cloud.ibm.com/docs/en/api/qiskit/synthesis},
accessed 7 November 2026.
\bibitem{fredkinshallow} P. Cruz and B. Murta, Shallow unitary decompositions
of quantum Fredkin and Toffoli gates for connectivity-aware equivalent circuit
averaging, APL Quantum 1, 016105 (2024),
\href{https://arxiv.org/abs/2305.18128}{arXiv:2305.18128}.
\bibitem{peterson} E. C. Peterson, G. E. Crooks, and R. S. Smith,
Fixed-depth two-qubit circuits and the monodromy polytope, Quantum 4, 247
(2020), \href{https://doi.org/10.22331/q-2020-03-26-247}{doi:10.22331/q-2020-03-26-247};
\href{https://arxiv.org/abs/1904.10541v4}{corrected arXiv version v4} (2021).
\bibitem{qubitization} G. H. Low and I. L. Chuang, Hamiltonian simulation by
qubitization, Quantum 3, 163 (2019),
\href{https://doi.org/10.22331/q-2019-07-12-163}{doi:10.22331/q-2019-07-12-163}.
\bibitem{cirqsyc} The Cirq Developers, Analytical decompositions into
Sycamore gates, Cirq source module \texttt{two\_qubit\_to\_sycamore.py},
\url{https://github.com/quantumlib/Cirq/blob/main/cirq-google/cirq_google/transformers/analytical_decompositions/two_qubit_to_sycamore.py},
accessed 7 November 2026; the audit uses Cirq 1.3.0.
\bibitem{hp24} K. Huang and J. Palsberg, Compiling conditional quantum gates
without using helper qubits, Proceedings of the ACM on Programming Languages
8, PLDI, Article 206 (2024),
\href{https://doi.org/10.1145/3656436}{doi:10.1145/3656436}.
\bibitem{vale24} R. Vale et al., Decomposition of multi-controlled special
unitary single-qubit gates,
\href{https://arxiv.org/abs/2302.06377}{arXiv:2302.06377} (2023).
\bibitem{tk2} Quantinuum, pytket transform API: KAKDecomposition,
\url{https://docs.quantinuum.com/tket/api-docs/transform.html}, accessed
7 November 2026. One unrestricted-angle TK2 represents any two-qubit Cartan class.
\end{thebibliography}
